\documentclass[11pt,a4paper]{article}
\usepackage[utf8]{inputenc}
\usepackage[T1]{fontenc}
\usepackage[french]{babel}
\usepackage{amsmath,amssymb}
\usepackage{geometry}
\usepackage{hyperref}
\geometry{margin=2.2cm}

\title{ADPro pour politiques de diffusion robotiques\\\large Explication des étapes et implémentation dans le code}
\author{Projet SIMUL -- Franka Panda}
\date{\today}

\begin{document}
\maketitle

\section{Objectif}
ADPro (\emph{Adaptive Diffusion Policy}) est une adaptation \emph{test-time} d'une politique de diffusion pré-entraînée.
L'idée est de ne pas réentraîner le modèle, mais de guider le débruitage avec de la géométrie de tâche.

Dans ce projet, ADPro est implémenté dans \texttt{models/adpro.py} sous la classe \texttt{ADPro}, avec deux variantes:
\begin{itemize}
\item \texttt{implementation='paper'}: version la plus fidèle au papier (FGR + guidance Chamfer + contrainte sphérique).
\item \texttt{implementation='practical'}: version plus robuste/pragmatique en simulation (target explicite + blend optionnel).
\end{itemize}

\section{Notation}
\begin{itemize}
\item Observation: $O = (O^0, O^1)$ avec $O^0$ nuage de points pince et $O^1$ nuage de points scène.
\item Action 7D: $a=[x,y,z,q_x,q_y,q_z,q_w]$.
\item Étape de diffusion: $t$ (du bruit fort vers faible).
\item Politique de base: réseau $\epsilon_\theta(O,a_t,t)$.
\end{itemize}

\section{Étape 1 -- Initialisation consciente de la tâche}
\subsection{Principe théorique}
Au lieu de partir d'un bruit gaussien pur $a_T\sim\mathcal N(0,I)$, ADPro calcule une pose initiale structurée $a_0^{\text{est}}$ via la géométrie de scène,
puis la projette au niveau de bruit $t=M$:
\begin{equation}
 a_M = \sqrt{\bar\alpha_M}\,a_0^{\text{est}} + \sqrt{1-\bar\alpha_M}\,\epsilon,\quad \epsilon\sim\mathcal N(0,I).
\end{equation}

\subsection{Implémentation}
Fonction: \texttt{\_task\_aware\_init(obs, t\_start)}.
\begin{itemize}
\item Sous-échantillonnage des points: \texttt{subsample\_pointcloud}.
\item Filtrage des points objet (pour réduire l'effet table): masque en $z$.
\item Variante \texttt{paper}: estimation rigide complète par \texttt{fast\_global\_registration} puis conversion via \texttt{se3\_to\_action}.
\item Variante \texttt{practical}: cible explicite de grasp (via \texttt{\_target\_grasp\_action}) + correction par centroïdes.
\end{itemize}

\section{Étape 2 -- Un pas de débruitage diffusion}
\subsection{Principe}
On exécute d'abord le pas DDPM standard pour obtenir $\tilde a_{t-1}$.

\subsection{Implémentation}
Dans \texttt{\_adpro\_step}:
\begin{itemize}
\item Prédiction du bruit: \texttt{eps\_pred = policy.network.forward(obs, a\_t, t)}.
\item Pas DDPM: \texttt{a\_tilde = policy.ddpm\_step(...)}.
\item Option déterministe: \texttt{deterministic\_denoising=True} coupe le bruit de reverse-process pour plus de stabilité.
\end{itemize}

\section{Étape 3 -- Estimation de l'action propre (formule type Tweedie)}
\begin{equation}
\hat a_0 = \frac{a_t - \sqrt{1-\bar\alpha_t}\,\epsilon_\theta(O,a_t,t)}{\sqrt{\bar\alpha_t}}.
\end{equation}
Implémentation: \texttt{\_predict\_clean\_action}.

\section{Étape 4 -- Guidance sur manifold de tâche}
\subsection{Version papier (surface-à-surface)}
On construit la perte géométrique sur nuages de points (Chamfer):
\begin{equation}
\mathcal L = d_{\text{Chamfer}}\big(T(\hat a_0)\,O^0,\;O^1\big).
\end{equation}
Puis on corrige l'action intermédiaire dans la direction opposée au gradient.

Implémentation dans \texttt{\_apply\_task\_manifold\_guidance}:
\begin{itemize}
\item Conversion action $\to$ pose: \texttt{action\_to\_se3}.
\item Remise en repère local pince puis application de $T(\hat a_0)$ (cohérence des repères).
\item Gradient Chamfer: \texttt{chamfer\_distance\_grad}.
\end{itemize}

\subsection{Version pratique (point-à-point)}
La guidance pousse vers une \texttt{target\_grasp\_action} explicite (position + quaternion cible) via une erreur L2 simplifiée.

\section{Étape 5 -- Contrainte sphérique (Eq. 9)}
La correction est normalisée puis mise à l'échelle par le bruit courant:
\begin{equation}
 a_{t-1} = \tilde a_{t-1} - s\,\sqrt{d}\,\sigma_t\,\frac{\nabla_a\mathcal L}{\|\nabla_a\mathcal L\|},
\end{equation}
où $d$ est la dimension d'action et $s$ est un facteur de stabilisation (dans le code: \texttt{spherical\_scale}).

Implémentation:
\begin{itemize}
\item \texttt{use\_spherical=True}: active la formule normalisée.
\item \texttt{spherical\_scale}: contrôle l'intensité du pas.
\end{itemize}

\section{Étape 6 -- Boucle complète d'inférence}
Dans \texttt{inference(...)}:
\begin{enumerate}
\item Choix de l'initialisation: task-aware (si \texttt{use\_fgr}) ou gaussienne.
\item Boucle décroissante en $t$ de \texttt{t\_start} vers $0$.
\item À chaque itération: DDPM $\to$ estimation $\hat a_0$ $\to$ guidance manifold $\to$ contrainte sphérique.
\item Sortie finale: action $a_0$ appliquée à l'environnement.
\end{enumerate}

\section{Paramètres importants et rôle}
\begin{itemize}
\item \texttt{implementation}: \texttt{paper} vs \texttt{practical}.
\item \texttt{M}: niveau de bruit initial (plus petit = init plus proche de l'action finale).
\item \texttt{spherical\_scale}: stabilité vs agressivité de la guidance.
\item \texttt{deterministic\_denoising}: stabilité temporelle en simulation.
\item \texttt{use\_blend} (practical): stabilisation additionnelle vers la cible explicite.
\end{itemize}

\section{Lien avec la visualisation}
Dans \texttt{simulate\_realtime.py}, l'option \texttt{--show-denoise} permet de visualiser:
\begin{itemize}
\item la trajectoire interne $a_t\to a_0$ (baseline vs ADPro),
\item l'initialisation ADPro,
\item les métriques de guidance (norme gradient, norme de pas, blend).
\end{itemize}

\section{Résumé}
ADPro améliore une politique de diffusion en injectant de la structure géométrique au test-time:
initialisation orientée tâche, guidance manifold, et normalisation sphérique des mises à jour.
La variante \texttt{paper} maximise la fidélité théorique; la variante \texttt{practical} maximise la robustesse en simulation.

\end{document}
